\subsection{Datasets}

\textbf{NADT-Prostate} (TCIA, DOI:10.7937/TCIA.JHQD-FR46): 39 patients, H\&E and ERG-stained
serial sections plus a PIN-4 cocktail, with real per-core Gleason score and phenotype
(benign/tumor) labels. Used to fit markers 1--3.
\textbf{PANDA} (Kaggle): 10{,}616 biopsy slides from two institutions (Karolinska, Radboud)
with official ISUP grade 0--5 labels; a stratified 1{,}200-slide subsample (1{,}137 after
tissue filtering) is used for external validation of markers 1--2.
\textbf{TCGA-PRAD} (GDC + cBioPortal \texttt{prad\_tcga\_pub}): 300 H\&E slides / 273 patients
with real Gleason score, PTEN copy-number, SPOP mutation, AR-activity score, and ERG fusion
status; also used, via a from-scratch relabeling of the GDC clinical API described below, for
recurrence-style outcome labels.
\textbf{SICAPv2}: used for diagnostic sanity checks (zero-shot text-prompt vs.\ linear-probe
comparison) reported in Section~\ref{sec:results}.
\textbf{LEOPARD} (\texttt{leopard.grand-challenge.org}): 508 patients, single institution
(Radboudumc), real biochemical-recurrence event and follow-up-time labels
(\texttt{case\_id}/\texttt{event}/\texttt{follow\_up\_years}; no clinical covariates in the
public release). Obtained from the dataset's public S3 bucket; subject to a publication
embargo pending the official challenge and baseline papers, which will be re-confirmed with
organizers before submission.
\textbf{PRECISE} (Zenodo, DOI:10.5281/zenodo.20721779): 25 patients, one core-needle biopsy
session each except one patient with three sessions (27 imaging sessions total, confirmed
against the dataset's own participant manifest and README, which also confirms these are
core-needle biopsies, not resections -- the manifest's age field is explicitly documented as
``age at biopsy''), paired H\&E and CKAPM+racemase IHC whole-slide images with 24{,}387
pixel-level expert annotations across seven classes (background, tumor, benign gland, artifact,
HGPIN, intraductal carcinoma, atypical intraductal proliferation, stroma) and, separately, real
per-image pathologist Gleason score.

\subsection{Foundation models and embedding extraction}

CONCH (\texttt{hf\_hub:MahmoodLab/conch}, ViT-B/16, 512-dim, \texttt{proj\_contrast=False})
and Virchow (\texttt{hf\_hub:paige-ai/Virchow}, ViT-H/14, 2560-dim = concatenated CLS token
and mean of patch tokens, requiring \texttt{mlp\_layer=SwiGLUPacked, act\_layer=SiLU}) are used
as frozen encoders; no weights are fine-tuned anywhere in this study. Tiles are 448$\times$448
(CONCH) or 224$\times$224 (Virchow). Whole-slide images without a usable multi-resolution
pyramid (LEOPARD, PRECISE) are tiled via windowed \texttt{zarr} reads against the
baseline-resolution page rather than full in-memory decoding.

Tile selection is \emph{not} uniform across cohorts; each cohort's actual method, with its
tissue-fraction threshold and per-slide tile count, is as follows. \textbf{NADT-Prostate and
PANDA} (markers 1--3 fitting and PANDA zero-shot transfer): random-uniform rejection sampling
-- a random $(y_0,x_0)$ tile origin is drawn repeatedly (up to 200 attempts per slide) at a
fixed pyramid level ($\approx 0.88\,\mu$m/px, chosen to match SICAPv2's native
$\approx 1\,\mu$m/px scale) and kept if the fraction of non-white pixels (mean grayscale
intensity $<205$) is $\geq 0.35$, until 16 tiles/slide are collected (PANDA deliberately reuses
NADT's exact tile count for a fair zero-shot comparison). \textbf{TCGA-PRAD} (markers 4--7,
ERG-fusion status, and all additional molecular candidates): the same random-uniform
rejection method and $0.35$ tissue-fraction threshold, but at a target $\approx 0.88\,\mu$m/px
($\approx 0.97\,\mu$m/px for the Virchow scale-matched re-embedding used in the marker-7 cross-
encoder check) derived \emph{dynamically per slide} from each file's AppMag tag rather than a
fixed pyramid-level index, after finding that a fixed-index fallback silently degraded every
TCGA-PRAD embedding to a non-tiled $\approx 17\,\mu$m/px thumbnail page whenever a slide's
resolution tag was unparseable (true for every slide in this cohort); 16 tiles/slide (64 for
the Virchow scale-matched re-embedding, matching LEOPARD Virchow's tile count for that specific
comparison). \textbf{LEOPARD}: non-overlapping grid sampling (not random) restricted to each
slide's own provided tissue-segmentation mask, kept if the mask-fraction within a tile is
$\geq 0.5$ (a real segmentation mask, not a non-white-pixel heuristic), randomly subsampled to
64 tiles/slide when the grid yields more candidates. \textbf{PRECISE}: non-overlapping grid
sampling against the dataset's pixel-level expert annotation mask, at a $150\,\mu$m physical
window (native $\approx 0.243\,\mu$m/px, resized to 448$\times$448) chosen specifically because
the standard $\approx 0.88\,\mu$m/px scale produced windows too coarse to find any qualifying
single-class region in this cohort's small annotated areas; a grid cell is kept if one
non-background class occupies $\geq 70\%$ of it, capped at 20 tiles per class per image (not
per slide overall, since the goal is balanced representation across tissue classes rather than
a fixed total). Per-slide embeddings are the mean of the tiles collected under each cohort's own
method above; per-patient values, where used, are the mean over a patient's slide-level
embeddings (or, for PRECISE, session-level embeddings -- see the PRECISE-specific note below).
Physical scale mismatch between training and application cohorts independently degrades
transfer regardless of sampling method (Section~\ref{sec:results}).

\subsection{Candidate markers}

Seven candidate markers are defined, each as a target and a training cohort: (1) H\&E
$\rightarrow$ Gleason score (NADT, RidgeCV); (2) H\&E $\rightarrow$ phenotype, tumor vs.\
benign (NADT, logistic regression); (3) ERG-stained image $\rightarrow$ Gleason score (NADT,
RidgeCV); (4) H\&E $\rightarrow$ PTEN copy-number loss (TCGA-PRAD, logistic regression); (5)
H\&E $\rightarrow$ SPOP mutation (TCGA-PRAD, logistic regression); (6) H\&E $\rightarrow$
AR-activity score (TCGA-PRAD, RidgeCV); (7) H\&E $\rightarrow$ biochemical-recurrence risk, a
\emph{post-hoc} marker discovered in the course of the LEOPARD experiment (Section
\ref{sec:leopard-methods}), fit directly against recurrence rather than against any of markers
1--6's targets. All continuous probes use
\texttt{StandardScaler + RidgeCV(alphas=logspace(-2,6,25))}; all binary probes use
\texttt{StandardScaler + LogisticRegression(C=1.0, class\_weight="balanced", max\_iter=2000)}
-- one fixed protocol across every marker and cohort, adopted after finding that per-test
hyperparameter search could flip a borderline result (marker 3's TCGA-PRAD ERG-fusion-status
replication moved from nominally significant to null once standardized).

\subsection{Cross-validation and statistical protocol}

All within-cohort evaluation uses patient-disjoint \texttt{GroupKFold(n\_splits=5)} (a
patient's slides never split across train/test), motivated by an early finding that our own
rule-based algorithm's apparent slide-level signal on NADT-Prostate was pseudo-replication that
disappeared at the patient level. Primary metrics are patient-level (mean prediction vs.\ mean
or majority true label per patient); slide-level metrics are reported as secondary. Across the
13 distinct marker-hypothesis tests conducted on real ground truth in this project (the 6
initial frozen-pool markers plus 7 additional TCGA-PRAD molecular candidates that were not
promoted to the pool) and the four nested/refit incremental-value tests introduced by the
confounder audit, Benjamini--Hochberg FDR correction is applied as one 17-test family;
external-cohort re-tests and cross-encoder replications of an already-tested hypothesis are
excluded from this family, since they confirm rather than newly test a hypothesis. Confidence
intervals are patient-ID-cluster bootstrap (2{,}000 resamples).

\subsection{Confounder audit}

For markers whose target is plausibly confounded by grade (4, 6, and -- retrospectively -- 7),
three models are compared: clinical-only, image-only, and combined. For markers 4/6, each
outer patient-disjoint fold constructs the meta-model's training image feature from inner-fold
OOF predictions, refits the image probe on the complete outer-training fold, and applies both
probe and meta-model to untouched outer-test patients. Marker 7's image score is fixed because
it was trained externally on LEOPARD; only clinical and combined Cox coefficients are fit in
each TCGA-PRAD outer fold. Primary statistics are held-out patient-level $\Delta$AUROC,
$\Delta R^2$, or $\Delta$C-index with 2{,}000 patient-bootstrap confidence intervals. The
earlier full-cohort likelihood-ratio analyses are retained as auxiliary, explicitly
in-sample results.

The nonparametric check permutes patient labels within rounded patient-mean Gleason strata and
reruns the complete nested pipeline, including image-probe refitting for markers 4/6, for
2{,}000 repetitions. For marker 7, paired event/time outcomes are permuted within grade
strata; its externally trained image score remains fixed while all outer-fold Cox models are
refit. Markers 1 and 3 are
excluded from this audit because their target is grade itself, making a ``does the image add
information beyond grade'' test circular.

For marker 7 specifically, the grade-only audit above is repeated against a richer,
fully-adjusted clinical baseline: grade, age, ordinal T-stage, PSA, and binary surgical margin
status, since grade alone leaves other well-known recurrence covariates unadjusted. Age and
T-stage (simplified to major stage 1--4, discarding substage letters) come from the GDC
\texttt{cases} endpoint. PSA and margin status (collapsed to positive = R1/R2 vs.\ negative =
R0, with RX treated as missing) come from cBioPortal's \texttt{prad\_tcga\_pub} patient-level
clinical data. Both sources were confirmed directly to carry these fields at the patient level
(they are absent from the sample-level clinical JSON used elsewhere in this project). Any
patient missing a covariate is excluded from this specific analysis (complete-case), and the
resulting cohort size is reported alongside the result
(Section~\ref{sec:res-marker7-clinical}) rather than imputed.

We additionally fit a sequential clinical hierarchy with held-out predictions: M0 = grade;
M1 = grade+age; M2 = grade+PSA+pT stage; M3 = M2+surgical margin; M4 = M3+tissue-source
site; and M5 = M4+marker7 risk. Numerical variables are standardized inside each training
fold, site is one-hot encoded with unseen test sites ignored, and a regularized Cox model
($\alpha=1$) is fit in five patient-disjoint outer folds. Each model reports its own
complete-case $n$, event count, and bootstrap C-index interval; M4 and M5 are compared on the
same 153 patients with a paired patient bootstrap.

\subsection{Reliability tiers}

Each initial marker is assigned one of five prospectively frozen tiers, from strongest to weakest evidence:
\emph{externally transportable} (a fixed, source-cohort-fitted probe applied with no retraining
reproduces the association in an independent cohort); \emph{cross-cohort replicable} (the
association reproduces after refitting in an independent cohort or a within-cohort held-out
site, but no zero-shot transport test exists); \emph{internally supported, externally
untested} (internally consistent across encoders and/or molecular assays, but no independent
cohort with the required label exists, confirmed by an exhaustive dataset search);
\emph{context-sensitive} (direction or magnitude is unstable across site, scale, encoder, or
assay); \emph{unsupported/null} (fails multiple-comparison correction or is null across
independent encoders).

\subsection{LEOPARD 3-pool design and domain-shift diagnostic}
\label{sec:leopard-methods}

Following the protocol-freeze discipline above, three evidence pools were prospectively fixed \emph{before}
running the LEOPARD experiment: a \emph{naive} pool (markers selected by within-cohort
significance alone, without cross-encoder or confounder gating: 1, 2, 4, 6), a
\emph{qualified} pool (markers passing every qualification gate: 1, 2, 4 -- marker 3 is
inapplicable to LEOPARD, which has no ERG-stained images), and an \emph{all-candidate} pool
that additionally includes the null and context-sensitive markers (5, 6). Each pool's markers
are scored zero-shot (source-cohort-fitted probes, no LEOPARD data used in fitting) and
combined via a 5-fold cross-validated Cox proportional-hazards model; pools are compared by
Harrell's C-index, time-dependent AUROC, integrated Brier score, calibration slope, and a
likelihood-ratio test against a covariate-free null model (LEOPARD's public labels carry no
clinical covariates, so an ``incremental value over clinical baseline'' claim is not testable
with this cohort and is not attempted).

When this zero-shot comparison returned uniformly sub-chance concordance for every pool, we
ran a separate, explicitly exploratory \emph{domain-shift diagnostic}: rather than transfer an
existing probe, a fresh model (PCA to 4--64 components, then a lightly regularized Cox model,
5-fold cross-validated) is fit directly on LEOPARD's own embeddings against LEOPARD's own
recurrence labels, to distinguish ``no recoverable recurrence signal in these embeddings'' from
``the specific zero-shot probes fail to transfer.'' The same diagnostic, and the same landmark
and grade-stratified confounder checks described in Section~\ref{sec:results}, are repeated
with a second, independently trained encoder (Virchow) and, separately, applied zero-shot to
TCGA-PRAD (after reconstructing recurrence-style labels from GDC's \texttt{follow\_ups}
records, since a pre-existing label file in the repository could not be traced to a
generating script and failed a spot check against the GDC API) as a genuine external-cohort
test of the resulting post-hoc marker (marker 7).

This reconstructed label is itself benchmarked against TCGA's standardized TCGA-CDR endpoints
(Liu et al.\ 2018), carried by cBioPortal study \texttt{prad\_tcga\_pan\_can\_atlas\_2018}:
event-status agreement (Cohen's $\kappa$) against \texttt{PFS\_STATUS} and \texttt{DFS\_STATUS},
time correlation (Spearman $\rho$) against \texttt{PFS\_MONTHS} and \texttt{DFS\_MONTHS}, and a
direct re-scoring of marker 7's zero-shot transfer using PFS and DFS themselves as the outcome,
in place of our reconstructed label, for the same 270-patient embedded cohort.
The generating script also exports a 500-case provenance table containing every raw follow-up
record, record-level exclusion reason, endpoint decision rule, and final label. A stricter
sensitivity endpoint requires a documented tumor-free follow-up before a later
\texttt{wt-with tumor} event; first/only with-tumor observations are excluded as
persistent/indeterminate disease rather than treated as recurrence or censoring. We report
3- and 5-year landmark AUROC, 1--5-year time-dependent AUROC, and correlation between risk and
censoring time for each available endpoint.

\subsection{QWK conversion}

To compare marker 1 against literature figures reported as F1 or quadratic weighted kappa
(QWK) rather than Spearman correlation, continuous predictions are converted to ordinal
categories via quantile binning matched to the number of true grade categories present in each
cohort (a transparent conversion that does not use true-label thresholds beyond the category
count), and QWK is computed against the true ordinal grade. This is applied uniformly across
all five evaluated cohorts.
